\documentclass[addpoints]{exam}
\usepackage[utf8x]{inputenc}
\usepackage[ngerman]{babel}
\usepackage{listings}
\usepackage[top=1.5cm,bottom=0.5cm,headsep=0.5cm,headheight=3cm,%
left=1.5cm,right=1.5cm]{geometry}

\usepackage[T1]{fontenc}
\usepackage{booktabs} % schöne Tabellen
\usepackage{graphicx}
\usepackage{amsmath,textcomp} %diverses
\usepackage{enumitem}
\usepackage{upquote}
\usepackage{csquotes}

\lstset{language=Python, tabsize=4, basicstyle=\footnotesize, showstringspaces=false, mathescape=true, upquote=true}
\lstset{literate=%
  {==}{{=\hspace{0.15em}=}}3
  {Ö}{{\"O}}1
  {Ä}{{\"A}}1
  {Ü}{{\"U}}1
  {ß}{{\ss}}1
  {ü}{{\"u}}1
  {ä}{{\"a}}1
  {ö}{{\"o}}1
}
\begin{document}
\pointpoints{Punkt}{Punkte}
\bonuspointpoints{Bonuspunkt}{Bonuspunkte}                     
\renewcommand{\solutiontitle}{\noindent\textbf{Lösung:}%
\enspace}

\pagestyle{headandfoot}
\runningheadrule

\newcommand{\klaubez}{Aufgaben zu Fehlern, Debugger und einfachen Algorithmen}
\firstpageheader{Informatik }{\klaubez} {\thepage /\numpages}
\runningheader{Informatik }{\klaubez} {\thepage /\numpages}
\newcommand{\pfad}{../aufgabentex}
%-------------------------------------------------------------------
% Steuerung ueber \mitloesungen (0 = ohne Loesungen, 1 = mit Loesungen).
% Direktes Kompilieren dieser Datei -> 01bitsUndbytes.pdf (ohne Loesungen).
% Fuer die Loesungsversion siehe 01bitsUndbytes_loes.tex -> 01bitsUndbytes_loes.pdf.
\providecommand{\mitloesungen}{0}
\ifnum\mitloesungen=1
  \printanswers
\fi
%-------------------------------------------------------------------


% \frage{datei}: bindet eine Aufgabe ein. Fragetext und Loesungsbox werden in eine Box
% gesetzt und damit nie ueber einen Seitenumbruch getrennt (kein manuelles \newpage noetig).
\newsavebox{\qbox}
\newcommand{\frage}[1]{\setbox\qbox\vbox{\hsize\linewidth\input{\pfad/#1}}\box\qbox}


\begin{questions}


% Aufgaben zu Fehlern, Debugger und einfachen Algorithmen
\frage{fehler_01}
\frage{fehler_02}
\frage{fehler_03}
\frage{logisch_01}
\frage{debugger_01}
\frage{listen_01}
\frage{listen_02}
\frage{listen_03}
\frage{suche_01}
\frage{suche_02}
\frage{suche_03}
\frage{anpassen_02}
\frage{bubble_01}
\frage{bubble_03}

% -------------------------------------------------
\end{questions}
\end{document}
